% example-switching.tex
% Theme switching example

\DocumentMetadata{}
\documentclass{ltx-talk}

% Mehrere Themes laden
\usepackage{lttalk-theme-magpie}
\usepackage{lttalk-theme-hawk}
\usepackage{lttalk-theme-sparrow}
\usepackage{lttalk-theme-anchovy}
\usepackage{lttalk-theme-minimal}
\ExplSyntaxOn
\cs_new_protected:Npn \example_progress_theme:
  {
    \lttalksetup{right-margin=true,margin-width=2.2cm}
    \definelttalkslot{example-progress}{right}{custom}
      { role = navigation }
    \setlttalkslotcontent{example-progress}
      {Start\\\textbf{Themen}\\Abschluss}
    \positionlttalkslot{example-progress}{center}{4cm}{2cm}{2.5cm}
    \stylelttalkslot{example-progress}
      {font=\scriptsize\sffamily,color=lttalk-accent,align=left}
  }
\RegisterLTXTalkTheme{progress}{\example_progress_theme:}
\ExplSyntaxOff
\uselttalktheme{magpie}
\addlttalktheme{progress}
\addlttalktheme{anchovy}

\title{Theme-Wechsel Demonstration}
\author{Max Mustermann}
\date{\today}

\begin{document}

% Start mit Magpie und dem Nur-Farb-Theme Anchovy
\begin{frame}
  \frametitle{Magpie Theme mit Anchovy-Farben}
  Diese Folie verwendet das von Madrid inspirierte Magpie-Theme.
  \begin{itemize}
    \item Header mit Titel und Navigation
    \item Linke Seitenleiste mit Logo
    \item Footer mit Institut und Datum
    \item Additives Teiltheme mit eigener Pfadnavigation
  \end{itemize}
\end{frame}

% Wechsel zu Hawk
\uselttalktheme{hawk}

\begin{frame}
  \frametitle{Hawk Theme}
  Diese Folie verwendet das von Hannover inspirierte Hawk-Theme.
  \begin{itemize}
    \item Linke Navigationsleiste
    \item Rechtsbündiger Folientitel
    \item Semantisch ausgezeichnete Navigation
  \end{itemize}
\end{frame}

% Wechsel zu Sparrow
\uselttalktheme{sparrow}

\begin{frame}
  \frametitle{Sparrow Theme}
  Diese Folie verwendet das von Szeged inspirierte Sparrow-Theme.
  \begin{itemize}
    \item Nur das Nötigste
    \item Dezente Linien
    \item Keine Ablenkungen
  \end{itemize}
\end{frame}

% Zurück zum Default
\uselttalktheme{default}

\begin{frame}
  \frametitle{Default ltx-talk}
  Diese Folie verwendet das Standard-Layout von ltx-talk.
  \begin{itemize}
    \item Kein Theme aktiv
    \item Originales ltx-talk-Verhalten
  \end{itemize}
\end{frame}

\end{document}
